\documentclass[tikz,border=6pt]{standalone}
\usepackage{peabra}
\begin{document}
\begin{tikzpicture}[x=1cm,y=1cm, a/.style={doboz=36mm, font=\footnotesize, minimum height=16mm}]
  \node[a] (f) at (0,0) {\textbf{hibaok} (fault)\\[2pt] például elállított szkennermező, szoftverhiba, kopott alkatrész};
  \node[kiemeltdoboz=36mm, font=\footnotesize, minimum height=16mm] (e) at (5,0) {\textbf{hiba} (error)\\[2pt] hibás belső állapot: a vezérlés „üresnek” látja a zónát};
  \node[draw=pePiros, fill=pePirosH, rounded corners=2pt, line width=1pt, text width=36mm, align=center, font=\footnotesize, minimum height=16mm] (h) at (10,0) {\textbf{hibajelenség} (failure)\\[2pt] a szolgáltatás eltér az elvárttól: nincs sebességcsökkentés};
  \draw[nyil] (f) -- node[above, felirat] {aktiválódás} (e);
  \draw[nyil] (e) -- node[above, felirat] {terjedés} (h);
  \draw[nyil, draw=peKiemel, dash pattern=on 3pt off 2pt] (h.south) -- ++(0,-0.9) -| node[felirat, pos=0.25, below, text=peKiemel] {a magasabb szint (a cella) számára ez hibaok} (f.south);
  \node[felirat, align=center, text width=150mm] at (5,-2.9) {véletlen hibaok: hardveröregedés, valószínűséggel jellemezhető; szisztematikus hibaok: specifikációs, tervezési, szoftveres hiba, minden azonos körülmények között újra jelentkezik};
  \node[felirat, align=left, anchor=west, text width=40mm, text=peKek] at (-1.8,1.65) {megelőzés és eltávolítás: tervezés, felülvizsgálat, teszt};
  \node[felirat, align=left, anchor=west, text width=46mm, text=peKek] at (7.2,1.65) {észlelés és tűrés: diagnosztika, redundancia, biztonságos állapot};
\end{tikzpicture}
\end{document}
